% Compile-only smoke test for legacy and lightweight interfaces.
\documentclass[UTF8,aspectratio=169]{beamer}
\usetheme[
  ColorDisplay=Blue,Miniframes=none,ContentMuticols=false,
  BIBMode=none,Background=false,SectionTOC=hide
]{academic}
\title{通用主题排版回归测试}
\author{Template CI}
\date{}
\begin{document}
\begin{frame}[t]{普通文字与极简强调}
  \ResearchLead{原生正文、标题和公式}
  这里是普通中文文字，与 \textbf{粗体强调} 形成层级。
  \ResearchHeading{研究对象}
  \[
    \mathbf{M}\ddot{\boldsymbol{\xi}}+
    \mathbf{C}\boldsymbol{\xi}=\mathbf{F}(t).
  \]
  \begin{ResearchNote}
    无图标、无卡片圆形编号的重点说明。
  \end{ResearchNote}
\end{frame}
\begin{frame}[t]{旧 ResearchCard 兼容}
  \begin{ResearchCard}[1]{整数编号应隐藏}
    旧接口正常接受；卡片仍可选用。
  \end{ResearchCard}
  \begin{ResearchCard}{未编号内容}
    内容可直接编辑。
  \end{ResearchCard}
\end{frame}
\begin{frame}[t]{简单自动网格}
  \begin{autolayout}[mode=balanced,layout=2]
    \begin{autoblock}{模型}
      质量与惯量。
    \end{autoblock}
    \begin{autoblock}{评价量}
      频率与运动响应。
    \end{autoblock}
  \end{autolayout}
\end{frame}
\end{document}
